
\subsubsection{2.1 Weight-Space Learning and Model Zoos}

Unterthiner et al. (2020) introduce ModelZooDataset and demonstrate that per-layer weight statistics (mean, variance, spectral norms --- the $\hat{W}_L$ encoder) achieve $R^2$ > 0.984 on CIFAR10-GS generalization prediction using gradient-boosted trees. This result establishes the reference performance ceiling against which the present work evaluates. $\hat{W}_L$ achieves high performance through approximately invariant statistics --- moment-based features that do not systematically encode channel order --- rather than through architectural invariance. Whether architectural invariance provides irreducible benefit over well-designed non-invariant statistics remains an open question that the present work investigates.

Eilertsen et al. (2020) extend weight-space learning to classification tasks, finding that raw weight footprints encode distinguishable signals about training conditions and optimizer choices. Schürholt et al. (2021, 2022) introduce self-supervised hyper-representations and model zoo benchmarks, establishing that weight populations contain rich transferable structure. These works motivate the need for encoders that faithfully represent the functional content of a network's weights.

\subsubsection{2.2 Architecturally Invariant and Equivariant Encoders}

The theoretical foundation for permutation-invariant set functions is established by Zaheer et al. (2017): any permutation-invariant function on a set can be decomposed as $\rho$($\Sigma \phi (x_{i}))$, making sum pooling the canonical invariant architecture. By construction, DeepSets sum pooling achieves OrbitVar = 0 exactly, up to floating-point precision.

Neural Functional Networks (NFN; Zhou et al., 2023) extend this to equivariant mappings over weight spaces, introducing NF-Layers (NPLinear) with parameter sharing tied to the CNN permutation group structure and HNPPool for invariant aggregation. NFN achieves Kendall's $\tau$ = 0.934 on CIFAR-10-GS generalization prediction. Zhou et al. report performance improvements but do not measure OrbitVar --- the within-orbit representational variance that the present work quantifies directly.

Navon et al. (2023) (DWSNet) derive the complete set of affine equivariant and invariant linear layers for deep weight spaces from symmetry group principles, formalizing that the correct symmetry group for CNNs requires coupled row-column permutations across adjacent layers. This coupled structure is a critical design constraint adopted in the present work; functional equivalence under permutation is verified before any encoding runs.

Kofinas et al. (2024) represent neural networks as parameter graphs and apply GNNs to learn equivariant embeddings, achieving state-of-the-art on several weight-space tasks. These approaches share the theoretical motivation for invariant encoding but do not directly measure how non-invariance translates to prediction error through the OrbitVar $\rightarrow$ MSE\_perm $\rightarrow R^2$ pathway.

\subsubsection{2.3 Permutation Symmetry in Weight Spaces}

The permutation symmetry of neural networks has been studied primarily in the context of loss landscape geometry (Entezari et al., 2022; Ainsworth et al., 2022) and neural network alignment \citep{Wang2020}. These works focus on finding the right permutation to align two networks (cross-model alignment), rather than on the within-model problem of measuring variance under all permutations of a single network.

Within-orbit variance (OrbitVar) and cross-model variance are orthogonal concerns: post-hoc alignment reduces the latter but cannot reduce the former, since OrbitVar is measured over all permutations of a single network, not across networks. The DWSNet formalism \citep{Navon2023Equivariant} provides the mathematical foundation for distinguishing these two problems; the MSE bias-variance decomposition over permutation orbits introduced here operationalizes this distinction into a directly measurable diagnostic.

\subsubsection{2.4 Position in the Literature}

The present work occupies a distinct position: rather than proposing a new encoder architecture, it provides the first empirical causal attribution of the encoder invariance $\rightarrow R^2$ relationship on a standard model zoo benchmark. The MSE decomposition framework (MSE\_res + MSE\_perm) is a reusable diagnostic applicable to any encoder to quantify its permutation sensitivity. The closest related measurement is the implicit comparison embedded in NFN's Kendall's $\tau$ improvement \citep{Zhou2023Permutation}, but that result conflates invariance benefits with representational capacity differences; the present controlled comparison --- CISE versus DeepSets at matched prediction capacity (same LightGBM predictor) --- isolates the invariance effect directly.

